\section{Related Work}

Initial research often prioritized specific aspects of UAV operations, such as path planning or flight dynamics. Swarmlab, which is a MATLAB based testbed designed for the rapid comparison of swarm algorithms which effectively visualizes metrics like swarm connectivity and safety. It relies on simplified point mass dynamics or basic quadcopter models, lacking the hardware level fidelity \cite{swarmlab}. On the other end tools such as RotorS provide high fidelity physics simulations within the Gazebo environment \cite{rotorS} capable of simulating sensors like IMUs and odometry and it mainly focusses on aerodynamic control on multi-UAVs rather than focussing on networking in swarm simulation. DronNS3 is built on the need for usability, a framework designed to lower the barrier to entry for UAV research. This emphasized and focused more on making ease of mission planning. It introduces a simplified python based autopilot library that abstracts MAVLink messages into high level commands \cite{drons3}. Although this provides and easy setup, visualization of the context could be identified as a drawback.

Simulation frameworks have been developed by integrating communication network. Integration of heterogeneous tools: which is a multi-layer architecture integrating Gazebo, ArduPilot and ns3 using a custom TCP based synchronization tool called GZUAVCHANNEL. This approach introduced latency and lacked support for real time telemetry data paths \cite{avens}. FlyNetSim bridges the gap between ArduPilot and ns3 without the overhead of virtual machines has been the main goal of introduction. The core innovation is the use of a middleware based on Zero Message Queue (ZMQ), which is a light weight message library. This intentionally avoids heavy weight physics engines like Gazebo to maintain scalability \cite{flynetsim}. GazeboNS3 could be identified as one of the recent advancements in integrating communication into the simulation. It proposes a Digital Twin system that combines Gazebo, ArduPilot and ns3. This system is capable of communicating dynamically using wireless mobile adhoc networks.

Edge processing (edge computing) is becoming increasingly popular in 
UAV networks which immensely save time and battery life. Studies have showed that use of advanced AI algorithms to plan UAV flight paths and decide where data should be processed to reduce delays and energy use \cite{edgeWorks}. However, these studies mostly focus on using UAVs as flying servers to process data for devices on the ground. They do not closely look at how a swarm of multiple UAVs handles processing its own data while flying together. Furthermore, existing research usually tests these ideas using basic mathematical simulations with simple, simulated data sizes. This leaves a major gap in understanding how heavy real world tasks like running computer vision models for human detection with realistic flight controllers like ArduPilot SITL.